\documentclass[a4paper]{article}

\usepackage[english]{babel}
%\usepackage[utf8x]{inputenc}
\usepackage{amsmath}
\usepackage{graphicx}
\usepackage[colorinlistoftodos]{todonotes}
\usepackage[all]{xy}
\usepackage{dsfont}
\usepackage{amssymb}
\usepackage{bbold}
\usepackage{amsthm}

\usepackage{tikz}
\usepackage{verbatim}
\usetikzlibrary{matrix}
\usetikzlibrary{cd}

\usepackage[autostyle=true]{csquotes}
\usepackage{pdfpages}
\usepackage{hyperref}
\usepackage{fancyhdr}


\newtheorem{theorem}{Theorem}[section]
\newtheorem{proposition}{Proposition}[theorem]
\newtheorem{lemma}{Lemma}[theorem]
\newtheorem{corollary}{Corollary}[theorem]
\newtheorem*{definition}{Definition}


\newcommand{\norm}[1]{\left\lVert#1\right\rVert}

\newcommand{\Rn}{\mathbb{R}^n}

\newcommand{\too}[1]{\overset{#1}{\longrightarrow}}
\newcommand{\hookar}[1]{\overset{#1}{\hookrightarrow}}
\newcommand{\colim}[1]{ \underset{#1}{\texttt{colim}} \text{ } }
\newcommand{\fset}[1]{ \{ #1_0 , \dots , #1_n \} }
\newcommand{\ddt}{\frac{d}{dt}}


\newcommand{\N}{\mathbb{N}}
\newcommand{\Z}{\mathbb{Z}}
\newcommand{\R}[1]{\mathbb{R}^{#1}}




\title{Taller 3}
\date{}
\fancyhead[R]{Universidad Finis Terrae \\ Facultad de Ingeniería \\ Algebra Lineal \\ 2do semestre de 2024}
\fancyhead[L]{\includegraphics[scale=0.1]{UFT.png}}

\begin{document}

\maketitle
\thispagestyle{fancy}
\begin{enumerate}


    \item Considere el conjunto $\{ ( 1 , \alpha , \alpha ), ( 1 , 2 , 1) , (\alpha , 2 , 1 )\}$.
    \begin{enumerate}
        \item   Determine para qu\'e valores de $\alpha$ el conjunto es li\-neal\-mente dependiente.
        \item 
 Determine la dimensión del espacio generado si $\alpha=2$.
    \end{enumerate}

\item Sea que $A$ es una matriz de orden $n$ y $x\in \mathbb{R}^n$ un vector no nulo. Supongamos
que $A^2x = 0$, pero $Ax \neq 0$. Pruebe que $x$ y $Ax$ son linealmente independientes.


    \item Sea $H$ el subconjunto de $\mathbb{R}^3$ definido por
\[H=\{(x,y,z)\in\mathbb{R}^3\colon 2x-3y+z=0  \}\]

\begin{enumerate}
    \item  Demuestre que $H$ es un subespacio vectorial de $\mathbb{R}^3$.
    \item  Encuentre una base de $H$ y determine su dimensión.
    \end{enumerate}
\item Considere el conjunto de vectores $S = \{( 1 , 0 , 0 ), ( 2 , 1 , 0 ) , ( 1 , 1 , 1 )\}$ en $\mathbb{R}^{3}$.
\begin{enumerate}
    \item Muestre que $S$ es base de $\mathbb{R}^{3}$.
    \item Determine las coordenadas de $ ( 2 , 1 , 1 )$ en esta base.
    \item Encuentre una base ortonormal de $\mathbb{R}^{3}$ a partir de $S$ usando el método de Gram-Schmidt.
    \item Encuentre la proyección de $ ( 2 , 1 , 1 )$ sobre el plano generado solo por los vectores $\{( 1 , 0 , 0 ), ( 2 , 1 , 0 )\}$.
\end{enumerate}

\end{document}